% ===================================================================
%  بخش ۲ — مبانی نظری و دامنه‌ی اعتبار توصیف همسانگرد  (کاملاً نو)
%  منبع: draft/section2.md  (کار ۳ در evidence/integration_plan.md)
% ===================================================================
\section{مبانی نظری و دامنه‌ی اعتبار توصیف همسانگرد}
\label{sec:theory}

\subsection{منشأ برچسب‌های کروی در توصیف گسیل و پراکندگی}
\label{sec:theory-origin}

چنان‌که در مقدمه اشاره شد، برچسب‌های $l$ و $m$ شناسه‌های جداسازی معادله‌ی
موج در دستگاه مختصات کروی‌اند:
%
\begin{equation}
  \nabla^2\psi + k^2\psi = 0
  \;\Longrightarrow\;
  \psi(r,\theta,\varphi) = R_l(r)\, Y_{lm}(\theta,\varphi).
  \label{eq:separation}
\end{equation}
%
اعتبار این برچسب‌گذاری به تقارن کروی مرز مسئله وابسته است
\cite{griffiths2018introduction, jackson1999classical, novotny2012principles}؛
اگر مرز تقارن کروی نداشته باشد، جداسازی متغیرها به این شکل انجام نمی‌شود و
$l$ دیگر عدد کوانتومی خوبی نخواهد بود.

باید دقت کرد که نانواپتیک مدرن خود از ابتدا توصیفی جهت‌دار از نرخ گسیل
به دست می‌دهد: نرخ واپاشی دوقطبی $\mathbf{p} = p\,\hat{\mathbf{n}}$ در
مکان $\mathbf{r}_0$ با چگالی موضعی \emph{جزئی} حالت‌ها،
%
\begin{equation}
  \frac{\Gamma}{\Gamma_0}
  = \frac{6\pi c}{\omega}\,
    \hat{\mathbf{n}}\cdot\operatorname{Im}\!\big[\overleftrightarrow{\mathbf{G}}(\mathbf{r}_0,\mathbf{r}_0;\omega)\big]\cdot\hat{\mathbf{n}},
  \label{eq:pldos}
\end{equation}
%
بیان می‌شود که در آن $\overleftrightarrow{\mathbf{G}}$ تانسور گرین دوگانه
است \cite[ch.~8]{novotny2012principles}. بنابراین پرسش این پژوهش
\emph{ابطال} ضرایب اینشتین نیست؛ پرسش این است که وقتی رابطه‌ی
\eqref{eq:pldos} را برای هندسه‌ای با محور ممتاز حل می‌کنیم، چه سهمی از
زبان کروی (برچسب‌های $l$ و $m$، توصیف تک‌پارامتری «فاکتور پورسل») هنوز
کارآمد می‌ماند و کجا باید کنار گذاشته شود. رابطه‌ی \eqref{eq:pldos}
هم‌زمان نشان می‌دهد که میانگین‌گیری روی جهت‌گیری دوقطبی —
$\tfrac{1}{3}\operatorname{Tr}\operatorname{Im}\overleftrightarrow{\mathbf{G}}$ —
تنها زمانی با نرخ واقعی برابر است که تانسور گرین همسانگرد باشد؛ در
هندسه‌ی نانومیله چنین نیست و همین نکته، مبنای مقایسه‌ی حالت‌های طولی و
عرضی در بخش~\ref{sec:results} است.

\subsection{تمایز اوربیتال اتمی از مد پلاسمونی}
\label{sec:theory-distinction}

پیش از ورود به بحث اصلی، تفکیک دو مفهومی که ریاضیات مشترکی دارند ضروری
است. برچسب‌های $l$ و $m$ هم در توصیف اوربیتال‌های اتمی و هم در توصیف مدهای
پلاسمونی یک نانوذره‌ی کروی به کار می‌روند، اما این دو از نظر فیزیکی یکسان
نیستند (جدول~\ref{tab:labels}).

\begin{table}[htbp]
  \centering
  \caption{مقایسه‌ی خاستگاه برچسب‌های $l$ و $m$ در دو زمینه‌ی متفاوت.}
  \label{tab:labels}
  \begin{tabular}{lll}
    \toprule
    ویژگی & اوربیتال اتمی & مد پلاسمونی نانوذره \\
    \midrule
    معادله‌ی حاکم & شرودینگر & ماکسول (هلمهولتز) \\
    ماهیت میدان  & تابع موج اسکالر $\psi$ & میدان \textbf{برداری} $\mathbf{E}$ \\
    منشأ تقارن   & پتانسیل کولنی $1/r$ & شکل مرز هندسی \\
    معنای $l$    & عدد کوانتومی مدار & برچسب مرتبه‌ی چندقطبی \\
    \bottomrule
  \end{tabular}
\end{table}

آنچه میان دو ستون جدول~\ref{tab:labels} مشترک است، ساختار زاویه‌ای
جواب‌ها است، نه فیزیک حاکم بر آن‌ها. در مورد اتم، تقارن کروی از پتانسیل
کولنی می‌آید؛ در مورد نانوذره، از شکل مرز هندسی. بنابراین نقد ارائه‌شده در
این پژوهش در سطح \textbf{دستگاه مختصات و تقارن مرزی} وارد می‌شود و به هیچ
روی ناظر بر درستی مکانیک کوانتومی یا نظریه‌ی می برای ذره‌ی کروی نیست.

این تمایز از آن رو اهمیت دارد که انتقال زبان اتمی به پلاسمونیک در برخی
مدل‌های پرکاربرد آگاهانه و صریح انجام شده است. مدل هیبریداسیون پلاسمونی،
پاسخ نانوساختارهای مرکب را با قیاس مستقیم از نظریه‌ی اوربیتال مولکولی و از
راه جفت‌شدگی مدهای چندقطبی با برچسب‌های $l$ و $m$ توصیف می‌کند
\cite{prodan2003hybridization}؛ همین قیاس در متن مرجع پلاسمونیک نیز
مبنای توصیف مدهای نانوپوسته‌ها قرار گرفته است
\cite[pp.~85--87]{maier2007plasmonics}. چنین قیاسی سودمند و در بسیاری موارد
پیش‌بین است، اما مانند هر قیاس دیگری دامنه‌ی اعتبار خود را دارد و این
دامنه به تقارن هندسه‌ی مورد بررسی گره خورده است. در همین راستا، ادبیات
نانوفوتونیک برای ساختارهای فلزیِ ناهمسانگرد به‌طور فزاینده از چارچوب
مفهومی \textbf{نانوآنتن} بهره می‌گیرد که در آن مشخصه‌های کلیدی، بازده
تابشی و الگوی تابش‌اند، نه مرتبه‌ی چندقطبی \cite{novotny2011antennas}.

\subsection{افت تقارن: چه چیزی در گذار از کره به استوانه باقی می‌ماند}
\label{sec:theory-symmetry}

گذار از یک نانوذره‌ی کروی به یک نانومیله، از دیدگاه تقارن، کاهش گروه تقارن
پیوسته‌ی $SO(3)$ به گروه $C_{\infty v}$ است: تقارن چرخشی حول
\textbf{همه‌ی} محورها جای خود را به تقارن چرخشی حول \textbf{یک} محور ممتاز
(محور طولی $z$) می‌دهد. در نتیجه، جداسازی متغیرها به شکل
رابطه‌ی~\eqref{eq:separation} دیگر برقرار نیست. پیامد این کاهش برای
برچسب‌گذاری مدها یکنواخت نیست
و باید با دقت تفکیک شود. تقارن چرخشی باقی‌مانده حول $z$ تضمین می‌کند که
عدد آزیموتالی $m$، که وابستگی میدان به مؤلفه‌ی $\varphi$ را به شکل
$e^{im\varphi}$ توصیف می‌کند، همچنان یک عدد کوانتومی خوب باقی بماند؛
همچنین در ساختار کشیده و شبه‌یک‌بعدی، ثابت انتشار محوری $\beta$ به عنوان
دومین شناسه‌ی مد ظاهر می‌شود. در برابر آن، عدد $l$ و تبهگنی
$(2l+1)$-گانه‌ی متناظر با آن، که هر دو نتیجه‌ی مستقیم تقارن کامل $SO(3)$
هستند، معنای خود را از دست می‌دهند. به همین دلیل است که در ادبیات
نانوسیم‌ها و نانومیله‌های فلزی، مدهای پلاسمونی معمولاً تنها با شماره‌ی
آزیموتالی برچسب می‌خورند و از برچسب چندقطبی کروی استفاده نمی‌شود؛ به عنوان
نمونه، در بررسی جفت‌شدگی نور با یک سیم طلای استوانه‌ای، بیشینه‌ی جفت‌شدگی
به پلاریتون پلاسمون سطحی از مرتبه‌ی $m = 2$ نسبت داده شده است
\cite{lee2008polarization, schmidt2008long}. نشانه‌ی تجربی این افت تقارن
در نانومیله‌ها آشکار است: طیف خاموشی نانومیله‌ی طلا به دو تشدید مجزای
طولی و عرضی شکافته می‌شود که پهنای خط و انرژی متفاوتی دارند
\cite{sonnichsen2002drastic}؛ در توصیف شبه‌استاتیکی نیز همین شکافت با
فاکتورهای هندسی بیضوی‌گون بازنمایی می‌شود
\cite[p.~72]{maier2007plasmonics}. چنین شکافتی در نانوذره‌ی کروی وجود ندارد و
مستقیماً پیامد وجود محور ممتاز است.

از این تحلیل، صورت‌بندی دقیق ادعای این پژوهش به دست می‌آید. آنچه در گذار
به تقارن استوانه‌ای رخ می‌دهد، بی‌اعتبار شدن یکجای برچسب‌گذاری کروی نیست،
بلکه حفظ بخشی از آن ($m$) و از دست رفتن بخش دیگر ($l$) است. تشخیص این
مرز، خود معیاری برای تعیین دامنه‌ی اعتبار مدل فراهم می‌کند. این
صورت‌بندی یک پیش‌بینی آزمون‌پذیر هم دارد: چون تنها $m$ حفظ می‌شود، پاسخ
نانومیله باید به جهت‌گیری دوقطبی نسبت به محور ممتاز به‌شدت وابسته باشد
و این وابستگی باید بسیار قوی‌تر از وابستگی جهتی در کره‌ی هم‌حجم باشد.
بخش~\ref{sec:results-orientation} این پیش‌بینی را با اجرای مستقیم دوقطبی
طولی و عرضی روی یک نوک می‌آزماید؛ آزمون کامل آن (دوقطبی مماسی روی کره)
در فهرست کارهای باقی‌مانده آمده است.

هندسه‌ی مورد بررسی در این پژوهش وضعیتی حتی پیچیده‌تر از یک استوانه‌ی
آرمانی دارد. ساختار از یک بدنه‌ی استوانه‌ای با دو کلاهک نیم‌کروی تشکیل شده
است؛ بنابراین مرز جسم در ناحیه‌ی بدنه با دستگاه مختصات استوانه‌ای و در
ناحیه‌ی کلاهک با دستگاه مختصات کروی هم‌خوان است، اما هیچ‌یک از این دو
دستگاه به تنهایی مرز \textbf{کل} جسم را توصیف نمی‌کند. در نتیجه، نه پایه‌ی
هارمونیک‌های کروی و نه پایه‌ی مدهای استوانه‌ای، هیچ‌کدام مجموعه‌ی طبیعی و
برازنده‌ای برای بسط میدان در این هندسه نیستند. همین ویژگی، که مستقل از هر
ملاحظه‌ی عددی دیگری برقرار است، توسل به حل مستقیم و تمام‌برداری معادلات
ماکسول را نه یک انتخاب محاسباتی بلکه یک ضرورت می‌سازد.

باید تأکید شود که هدف این بخش تنها تعیین حدود اعتبار برچسب‌گذاری‌های مبتنی
بر تقارن کروی است. توسعه‌ی یک پایه‌ی مدی متناسب با تقارن استوانه‌ای و
ارزیابی کارایی آن، خارج از دامنه‌ی این پژوهش است.
